\documentclass[11pt,a4paper]{article}
\usepackage[a4paper,margin=13mm,top=11mm,bottom=11mm]{geometry}
\usepackage[T1]{fontenc}
\usepackage{mathptmx}
\usepackage{booktabs,tabularx,array,enumitem,xcolor}
\usepackage[colorlinks=true,urlcolor=blue]{hyperref}
\pagestyle{empty}
\setlength{\parindent}{0pt}\setlength{\parskip}{2pt}
\setlist{nosep,leftmargin=12pt,topsep=1pt}
\newcolumntype{L}{>{\raggedright\arraybackslash}X}
\newcolumntype{P}[1]{>{\raggedright\arraybackslash}p{#1}}
\definecolor{ok}{HTML}{1f7a45}\definecolor{warn}{HTML}{a15c00}
\newcommand{\sect}[1]{\vspace{3pt}{\large\bfseries #1}\par\vspace{1pt}}
\begin{document}
{\LARGE\bfseries SeizeLUT}\hfill{\small BISP project, Dept.\ of ECE, NIT Warangal, November 2026}\par
{\large A self-testing, multiplier-free seizure detector on an FPGA, with doctor alerts}\par
{\small Akshara Bonam (23ECB0A11) \quad Sai Chaithr (23ECB0A20) \quad Sowjanya (23ECB0A22)}\par
\vspace{2pt}\hrule\vspace{3pt}

\textbf{In one line:} behind-the-ear EEG, ECG, EMG and motion go into an FPGA that computes the features and
classifies every second with 50 lookup tables (no multiplier, 1,202 LUTs in all), tests its own hardware, raises an
alarm LED and buzzer, and an ESP32 tells the doctor what to do.

\sect{What is done (all numbers measured)}
\begin{tabularx}{\textwidth}{@{}P{38mm}L@{}}
\toprule
\textbf{Detection} & SeizeIT2, 125 patients, 11,626\,h; design fixed on 96 training patients, scored once on 29
held-out patients: event-F1 0.248, \textbf{0.8 false alarms per day}; above TreeLUT, logic-gate network, INT8 network and
ChronoNet; two-stage version 73.3\,\% of seizures at 10.9 FA/h vs 71.1\,\% at 11.0 for the dataset paper's SVM\\
\textbf{Hardware cost} & classifier 398 LUTs, 468\,MHz, 0 DSP (6.7$\times$ fewer LUTs than TreeLUT, 13.2$\times$ fewer than a
logic-gate network); with the front end in the fabric \textbf{1,202 LUTs, no multiplier}, real time below 1\,MHz\\
\textbf{Self-testing (novelty)} & all 3,592 single hardware faults analysed on 8,714\,h; 255 selected test vectors detect
\textbf{every fault that changes a decision} (a normal test: 50.7\,\%); self-test takes 10.3\,\textmu s\\
\textbf{Board (Kria KV260)} & \textcolor{ok}{\textbf{all 9 board tests pass}}: 431,960 windows with 0 mismatches, all 3,592
injected faults as predicted, front end + classifier in the fabric equal to software on 9,214 windows; engine cycles
on the board equal simulation exactly; timing met at 100\,MHz; 3.2\,\% of the chip\\
\textbf{Demo} & live dashboard; red alarm LED and buzzer driven by the FPGA (working on the board)\\
\textbf{Doctor alert (ESP32)} & iPhone push + live page: one instruction (e.g.\ GO TO BED 3 NOW), heart / shaking /
muscle / brain tiles, escalation at 2 and 5 min, ``On my way'', call button. \textcolor{warn}{Tested in emulation; ESP32
hardware test pending}\\
\textbf{Paper} & 6-page IEEE paper with prior-work comparison and board results, ready\\
\bottomrule
\end{tabularx}

\sect{How good is it}
\begin{tabularx}{\textwidth}{@{}P{38mm}L@{}}
\textbf{Novelty: strong} & to our knowledge the first LUT-native seizure detector and the first complete fault analysis of a seizure
classifier, with a self-test proven on the board; the whole detector as multiplier-free hardware\\
\textbf{Evidence: strong} & strict patient split with sealed test patients; real board measurements, not only
simulation; a published baseline reproduced, including a protocol flaw in it\\
\textbf{Honest weaknesses} & LightGBM is more accurate (0.308) but needs about 30,000 LUTs; focal seizures still mostly
missed (as in the clinical literature); one retrospective dataset; doctor alert not yet tested on the ESP32\\
\textbf{Overall} & a complete system from data to working hardware, publishable at a good IEEE biomedical-circuits
conference; strong material for master's applications\\
\end{tabularx}

\sect{Where to submit (deadlines from the official pages, checked 20 Nov 2026)}
\begin{tabularx}{\textwidth}{@{}P{34mm}P{40mm}L@{}}
\toprule
\textbf{Conference} & \textbf{Deadline} & \textbf{Why / what to do}\\
\midrule
\textbf{1. IEEE EMBC 2027} (first choice) & full paper \textbf{24 Jan 2027}; decision 16 Apr & the IEEE EMBS flagship
conference; wearable detection fits; \textbf{new limit 4 pages + 1 page of references}: shorten the 6-page paper\\
\textbf{2. IEEE BioCAS 2027} (if EMBC rejects) & abstract \textbf{7 May 2027}; Oslo, 14--16 Oct & biomedical circuits and
systems: best fit for the whole system; comes after the EMBC decision, so no double submission\\
IEEE/ACM FCCM 2027 (hardware option) & abstract 10 Jan, paper 17 Jan 2027 & FPGA-focused; would compete with EMBC on
dates, so choose one of the two\\
\bottomrule
\end{tabularx}
{\small Never submit the same paper to two conferences at the same time.}

\sect{Next steps}
\begin{enumerate}[leftmargin=14pt]
\item Test the ESP32 alert with a real iPhone; film the full demo (alarm LED, buzzer, phone notification).
\item Shorten the paper to 4 pages + references for EMBC; ask the guide to review; submit before 24 January 2027.
\item Fill the guide's name, the certificates and signatures in the BISP report.
\end{enumerate}
\end{document}
